\documentclass[12pt,a4paper]{article}

% --- Paquetes de diseño y fuentes ---
\usepackage[utf8]{inputenc}
\usepackage[spanish, es-tabla]{babel}
\usepackage{amsmath, amsfonts, amssymb, amsthm}
\usepackage{graphicx}
\usepackage{geometry}
\usepackage{fancyhdr}
\usepackage{setspace}
\usepackage{xcolor}
\usepackage[colorlinks=true, linkcolor=blue!70!black, urlcolor=black, citecolor=blue]{hyperref}
\usepackage{tikz}
\usetikzlibrary{shapes.geometric, arrows, positioning}

% Configuración de márgenes
\geometry{top=2.5cm, bottom=2.5cm, left=2.5cm, right=2.5cm}
\onehalfspacing

% Cabecera y Pie de página
\pagestyle{fancy}
\fancyhf{}
\rhead{\small Semillero de Aprendizaje por Refuerzo}
\lhead{\small IA para Trading y Sentimiento}
\cfoot{\thepage}
\renewcommand{\headrulewidth}{0.4pt}

% Título y Autores
\title{
    \vspace{-1cm}
    \textbf{\Large IA y Emociones: Desarrollo de un Agente de Aprendizaje por Refuerzo para el Mercado de Criptomonedas} \\
    \vspace{0.5cm}
    \large Universidad Externado de Colombia \\
    \text{Ciencia de Datos -- Departamento de Matemáticas}
}
\author{
    \textbf{Julian Duarte} \quad \textbf{Julian Jimenez} \\
    \medskip
    \small \textit{Tutores:} \\
    \small Prof. Julián Fernando Sánchez \quad Prof. Juan Carlos Urueña \quad Prof. Camilo Castillo
}
\date{Abril de 2026}

\begin{document}

\maketitle

\tableofcontents
\newpage

\section{Introducción: ¿Por qué una IA que sienta el mercado?}

El mundo de las \textbf{criptomonedas} es conocido por ser como una montaña rusa. Muchas veces, estos cambios no ocurren solo por números, sino por el factor humano. Tradicionalmente, los programas de \textbf{trading} solo miran estadísticas, pero olvidan el sentimiento que mueve el mercado.

Este proyecto propone crear una Inteligencia Artificial (IA) que integre la psicología del mercado mediante el \textbf{Aprendizaje por Refuerzo}, entrenando al agente para reaccionar ante el pánico o la euforia colectiva.

Nuestra IA operará con la moneda de \textbf{Binance} (\textbf{BNB}) y un índice de Miedo y Codicia. El objetivo es que la IA aprenda a ser precavida cuando hay pánico y estratégica cuando hay euforia, buscando un crecimiento seguro del capital.

\section{¿Qué es el Aprendizaje por Refuerzo (RL)?}

Para lograr este objetivo de proteger el dinero y crecer de forma inteligente, necesitamos una herramienta que no solo siga reglas fijas, sino que aprenda de sus propios aciertos y errores. Aquí es donde entra en juego el Aprendizaje por Refuerzo (o \textit{Reinforcement Learning}), la tecnología que actúa como el ``sistema de aprendizaje'' de nuestra IA.

A diferencia de otros tipos de inteligencia artificial que aprenden mirando fotos o ejemplos estáticos, el RL aprende de la \textbf{experiencia directa}. Es muy parecido a cómo entrenamos a una mascota: si hace algo bien, recibe un premio; si hace algo mal, no recibe nada o recibe un ``castigo''.

Para que este proceso ocurra, siempre deben existir cinco piezas fundamentales que interactúan entre sí en un ciclo infinito:

\begin{enumerate}
    \item \textbf{El Agente}\textbf{:} Es el ``cerebro'' de la IA, el que toma las decisiones.
    \item \textbf{El Entorno:} Es el mundo donde vive el agente (en nuestro caso, el mercado de criptomonedas).
    \item \textbf{El Estado ($S_t$):} Es lo que el agente ve en un momento dado (su tablero de instrumentos).
    \item \textbf{La Acción}\textbf{:} Es la decisión que toma el agente (Comprar, Vender o Esperar).
    \item \textbf{La Recompensa ($R_t$):} Es el premio o castigo que recibe el agente por su acción.
\end{enumerate}

A partir de ahora, exploraremos una a una estas piezas para entender cómo nuestra IA se convirtió en una experta del mercado bajo un modelo donde el futuro solo depende del estado actual; esto es conocido como \textbf{Proceso de Decisión de Markov (MDP)}.

\section{El Agente: Nuestro Trader de Inteligencia Artificial}

En el centro de nuestro proyecto se encuentra el Agente. Si el mercado es el océano, el Agente es el navegante que debe aprender a manejar su barco de inversión. 

A diferencia de un trader humano que puede dejarse llevar por el pánico, nuestro Agente es un algoritmo puramente racional que toma decisiones basadas en la experiencia. Su única misión es aprender una \textbf{política óptima}\footnote{Es el conjunto de reglas o ``manual de conducta'' definitivo que el agente desarrolla para saber exactamente qué acción tomar en cada situación posible para ganar.} que le permita maximizar lo que en RL llamamos el \textbf{Return} o \textbf{Beneficio Acumulado}\footnote{Es la suma total de todas las recompensas acumuladas por el agente a lo largo de su vida, siendo el valor final que la IA intenta maximizar.} ($G_t$):

\begin{equation}
    G_t = R_{t+1} + \gamma R_{t+2} + \gamma^2 R_{t+3} + \dots
\end{equation}

Esta fórmula, aunque parezca compleja, tiene una lógica muy sencilla:
\begin{itemize}
    \item \textbf{$G_t$ (Meta Total):} Es la suma de todos los premios desde ahora hasta el futuro. ¿Por qué los sumamos? Porque el trading es una carrera de fondo, no una carrera de 100 metros. Sumar los premios permite que la IA sea estratégica: a veces, aceptará perder un poquito hoy si sabe que eso la llevará a una ganancia gigante más adelante; sin esta suma, con una mentalidad a corto plazo, la IA nunca arriesgaría nada para ganar más después.
    \item \textbf{$R$ (Recompensa):} Es el valor del premio o castigo que la IA recibe en un momento específico.
    \item \textbf{Los numeritos abajo ($t+1, t+2 \dots$):} Representan el tiempo: la próxima hora, la siguiente, y así sucesivamente. Nos dicen en qué momento exacto del futuro se recibirá cada premio.
    \item \textbf{$\gamma$ (El Factor de Paciencia):} Es un número entre 0 y 1 que decide qué tanto nos importa el futuro. 
    \begin{itemize}
        \item \textbf{¿Qué es?} Técnicamente se llama \textbf{factor de descuento}. Imagínalo como un control de ``visión estratégica'': si $\gamma$ es cercano a 1, la IA tiene una visión de largo alcance y valora mucho los beneficios que llegarán en el futuro. Si es cercano a 0, la IA sufre de una ``miopía estratégica'' y solo le interesa ganar dinero en la próxima hora, ignorando lo que pase después.
        \item \textbf{¿Por qué se eleva al cuadrado o al cubo?} Se eleva según el tiempo que falta ($\gamma^2, \gamma^3 \dots$) porque \textbf{un premio entre más lejos esté, menos valioso es hoy}. Es como si el premio se fuera ``enfriando'' con el tiempo. Esto obliga a la IA a ser eficiente: le dice que el futuro importa, pero que no puede quedarse esperando mil años a que pase algo bueno; tiene que actuar para ganar lo más pronto posible.
    \end{itemize}
\end{itemize}

\section{El Entorno: ¿Cómo ve la IA el mundo del Mercado?}

Para que nuestra IA aprenda a navegar en las finanzas, necesita un conjunto de información técnica y emocional conocido como el \textbf{Entorno de Observación}\footnote{Conjunto de variables que definen el mundo observable por la IA.}. Hemos construido este entorno sobre dos pilares:

\begin{itemize}
    \item \textbf{El Pilar Técnico (Binance):} Operamos \textbf{BNB} frente a \textbf{USDT}\footnote{Moneda estable (1:1 con dólar) utilizada como unidad de medida de rentabilidad.}. Elegimos la temporalidad de \textbf{1 hora} mediante la \textbf{vela horaria}\footnote{Representación gráfica del precio (apertura, cierre, máximos y mínimos) en intervalos de 60 minutos.} porque representa el ``punto de equilibrio'' estratégico: una temporalidad menor (como 1 o 5 minutos) saturaría a nuestra IA con \textit{ruido blanco} ---fluctuaciones insignificantes que no reflejan sentimientos reales---, mientras que una periodicidad diaria nos haría ignorar cambios críticos que en cripto ocurren en cuestión de horas. La hora es la medida ideal para que el agente reconozca tendencias emocionales sólidas y actúe con la agilidad necesaria.
    
    \item \textbf{El Pilar de Datos:} La \textbf{API de Binance}\footnote{Es el ``puente'' de comunicación que permite que nuestro programa se conecte con el mercado para obtener datos y enviar órdenes de forma automática.} nos provee datos con alta \textbf{liquidez}\footnote{Es la facilidad para comprar o vender un activo rápido y al precio actual, sin que nuestra operación mueva el mercado de forma brusca.}, garantizando que las decisiones de la IA sean ejecutables, considerando siempre la \textbf{volatilidad}\footnote{Medida de variación del precio; indica riesgo y oportunidad.} del mercado.

    \item \textbf{El Pilar Psicológico:} Utilizamos los datos de \textbf{Alternative.me}, una plataforma líder en análisis de sentimiento que procesa el pulso emocional del mercado cripto. Integramos su \textbf{Crypto Fear \& Greed}\footnote{Indicador psicométrico (0-100) que cuantifica el sentimiento emocional del mercado.} Index, el cual sintetiza diariamente múltiples fuentes de información: volatilidad técnica, impulso de volumen, actividad en redes sociales (como Twitter y Reddit), el dominio de Bitcoin y tendencias de búsqueda en Google. Esta integración permite que el agente no solo reaccione a precios, sino que comprenda la psicología colectiva que precede a los grandes movimientos del mercado.
\end{itemize}

El índice se mide en una escala de 0 a 100 siguiendo la lógica de las \textbf{Finanzas Conductuales}\footnote{Disciplina que estudia cómo los sesgos emocionales afectan el comportamiento de los mercados.}:
\begin{itemize}
    \item \textbf{0 - 24 (Miedo Extremo):} El pánico domina el mercado. Técnicamente, es el momento de compra ideal porque los precios suelen estar ``en oferta''.
    \item \textbf{25 - 49 (Miedo):} Cautela y desconfianza. El Agente suele esperar o comprar con cuidado.
    \item \textbf{50 - 74 (Codicia):} Optimismo. El mercado tiene fuerza y la gente sigue comprando.
\item \textbf{75 - 100 (Codicia Extrema):} Euforia peligrosa. Es el momento de venta o de proteger ganancias, ya que una caída suele estar cerca.
\end{itemize}

\subsection{El Poder del NLP frente a lo Determinista}

Una pregunta clave es: ¿por qué usar \textbf{Procesamiento de Lenguaje Natural}\footnote{Rama de la inteligencia artificial que se encarga de procesar y derivar significado del lenguaje humano para que las máquinas puedan entender el contexto y sentimiento tras las palabras.} \textbf{(NLP)} para el sentimiento en lugar de un indicador matemático tradicional? La respuesta reside en la naturaleza del mercado. Las métricas de análisis técnico convencional ---como las medias móviles--- son por naturaleza \textbf{deterministas}\footnote{Se refiere a modelos que arrojan siempre el mismo resultado para una entrada dada según reglas fijas, ignorando variables externas o cambios en el comportamiento humano.}: dependen de fórmulas matemáticas rígidas que solo procesan datos históricos, sin capacidad para interpretar el contexto emocional del momento.

El NLP, por el contrario, nos permite capturar la \textbf{causalidad emocional}. Al analizar el lenguaje en redes sociales y noticias, nuestra fuente (Alternative.me) detecta la intención y el miedo de los inversores antes de que estos ejecuten sus órdenes. Integrar este proceso no determinista en el estado del agente le otorga una "ventaja competitiva" psicológica, permitiéndole entender no solo \textit{qué} está pasando, sino \textit{por qué} está pasando.

Finalmente, nuestro laboratorio de pruebas abarca desde enero de 2023 hasta diciembre de 2024. \textbf{¿Por qué elegimos este periodo?} Tras considerar periodos más largos, seleccionamos este bienio porque captura la realidad actual del mercado: 2023 nos ofrece un escenario de recuperación y lateralidad tras la crisis de 2022, mientras que 2024 introduce el inicio de \textbf{ciclos alcistas}\footnote{También conocido como \textit{Bull Market}, es un periodo en el que los precios del mercado suben de forma sostenida en el tiempo, impulsados por el optimismo y la confianza de los inversores.} masivos. Este contraste es vital para que la IA aprenda a sobrevivir en la incertidumbre y a capitalizar en la euforia, evitando que se vuelva complaciente.

\section{El Estado ($S_t$): El Tablero de Instrumentos de la IA}

Si el mercado es un océano inmenso, nuestra IA no puede verlo todo a la vez. Para tomar decisiones, necesita un resumen de lo que está pasando exactamente en este momento. Este resumen se llama técnicamente \textbf{El Estado}\footnote{Conjunto de información disponible para el agente en un instante dado (precio, RSI, sentimiento, etc.).} \textbf{($S_t$)}.

Imagina que el Agente es el piloto de un avión. No necesita asomarse por la ventana para saber cómo volar; le basta con mirar su tablero de instrumentos. En nuestro proyecto, ese tablero tiene cuatro indicadores fundamentales que forman un vector matemático:

\begin{equation}
    S_t = \{M_t, RSI_t, Sent_t, I_t\}
\end{equation}

A continuación, explicamos qué mide cada uno de estos ``relojes'':

\begin{enumerate}
    \item \textbf{Régimen de Mercado ($M_t$):} Es el indicador de dirección. Nos dice si el precio en la última hora subió o bajó. Es como saber si el viento sopla a favor o en contra de nuestra posición.
    \item \textbf{El Termómetro Emocional ($Sent_t$):} Representa el \textbf{``momento psicológico''} del mercado. No es solo un número; es la captura de la euforia o el pánico colectivo en el instante $t$. Al incluirlo en el estado, le permitimos a la IA entender si el precio actual está influenciado por fundamentos económicos o por impulsos emocionales de los inversores.
    \item \textbf{El Medidor de Cansancio ($RSI_t$):} El \textbf{RSI}\footnote{Indicador de análisis técnico que mide la velocidad y el cambio de los movimientos de precios para identificar condiciones de sobrecompra o sobreventa.} es como el sensor de fatiga del precio. \textbf{¿Por qué el RSI?} Porque nos da una medida matemática de cuándo un movimiento de precio se ha extendido demasiado, permitiendo a la IA detectar ``burbujas'' locales o pánicos temporales que no son sostenibles, actuando como un contrapeso objetivo al índice de sentimiento.
    \item \textbf{El Interruptor de Inventario ($I_t$):} Es la memoria física de la IA. Le indica si en este momento tiene dólares en su bolsillo o si ya los gastó comprando BNB. 
\end{enumerate}

\textbf{¿Por qué simplificamos los datos? (Discretización):}
Para que el cerebro de nuestra IA sea eficiente, utilizamos una técnica llamada \textbf{Discretización}\footnote{Proceso de convertir datos continuos en categorías finitas para simplificar el aprendizaje del agente.}. 

En el mundo real, los datos son continuos (el sentimiento puede ser 42.5 o 42.6). Si la IA intentara aprender cada decimal, necesitaría una ``tabla de sabiduría'' infinita. Para evitar esto, convertimos los números en ``cajas'' estratégicas:
\begin{itemize}
    \item \textbf{Sentimiento ($Sent_t$):} El índice de 0 a 100 se agrupa en las 4 zonas que definimos (Miedo Extremo, Miedo, Codicia, Codicia Extrema). Así, para la IA es lo mismo un sentimiento de 10 que uno de 15; ambos significan: \textit{``La gente está aterrorizada, es probable que haya una oportunidad de compra''.}
    \item \textbf{RSI ($RSI_t$):} En lugar de ver 72.4, la IA solo ve la categoría: \textbf{``Mercado Agotado (Sobrecompra)''}.
\end{itemize}
Esto reduce drásticamente la complejidad del problema. Permite que la IA rellene su tabla de aprendizaje mucho más rápido, ya que solo tiene que aprender qué hacer en unas pocas situaciones clave (como ``Miedo alto + Precio bajo + Tengo dinero'') en lugar de perderse en un mar de decimales sin sentido estratégico.

\section{Arquitectura del Sistema: El Ciclo de Decisión}

Para facilitar la comprensión global, hemos diseñado el siguiente flujo operativo que resume desde la captura de datos hasta la generación de ganancias:

\begin{figure}[h]
    \centering
    \begin{tikzpicture}[node distance=1.8cm, auto]
        % Definición de estilos
        \tikzstyle{startstop} = [rectangle, rounded corners, minimum width=3cm, minimum height=1cm,text centered, draw=black, fill=blue!10]
        \tikzstyle{process} = [rectangle, minimum width=3cm, minimum height=1cm, text centered, draw=black, fill=orange!10]
        \tikzstyle{decision} = [diamond, aspect=2, minimum width=3.5cm, minimum height=1cm, text centered, draw=black, fill=green!10]
        \tikzstyle{arrow} = [thick,->,>=stealth]

        % Nodos del proceso con posicionamiento mejorado
        \node (data) [startstop] {Fuentes de Datos (Binance/Sentiment)};
        \node (proc) [process, below=of data] {Procesamiento y Discretización};
        \node (state) [process, below=of proc] {Definición del Estado ($S_t$)};
        \node (agent) [decision, below=1cm of state] {Agente RL (Q-Learning)};
        \node (market) [startstop, right=2.5cm of agent] {Entorno de Mercado};
        \node (eval) [process, below=1.5cm of agent] {Evaluación de Resultados};

        % Flechas con etiquetas limpias
        \draw [arrow] (data) -- (proc);
        \draw [arrow] (proc) -- (state);
        \draw [arrow] (state) -- (agent);
        \draw [arrow] (agent) -- node[above] {Acción} (market);
        \draw [arrow] (market.south) -- ++(0,-0.5) -| node[pos=0.25, below] {Recompensa} (agent.east);
        \draw [arrow] (agent) -- (eval);

    \end{tikzpicture}
    \caption{Flujo Operativo del Agente de Trading Emocional.}
\end{figure}

Este ciclo funciona de forma horaria: la IA capta la ``foto'' del mercado, decide qué hacer, recibe su premio o castigo y vuelve a empezar, acumulando sabiduría en cada interacción.

\section{Acciones y Recompensas: El Lenguaje del Éxito}

Para que el aprendizaje ocurra, la IA necesita saber qué puede hacer y cómo será evaluada. En el mundo de la IA, esto se define mediante las Acciones (nuestros movimientos) y la Recompensa (nuestra brújula de éxito).

\subsection{Nuestras Decisiones (Acciones $A_t$)}

El agente tiene tres movimientos posibles en cada hora de mercado. Imagínalos como los controles de un mando de consola:

\begin{itemize}
    \item \textbf{Mantener / Esperar ($A=0$):} Es la acción de ``no hacer nada''. Si tiene dólares, se queda en dólares; si tiene BNB, los conserva. Esta acción es vital porque enseña a la IA que, a veces, el mejor movimiento es quedarse quieto y observar.
    \item \textbf{Comprar ($A=1$):} La IA decide entrar al mercado. Cambia sus USDT por BNB. Solo puede hacerlo si su \textbf{inventario}\footnote{Estado de posesión del agente (posicionado con activos o en efectivo).} está vacío.
    \item \textbf{Vender ($A=2$):} La IA decide salir del mercado y transformar sus criptomonedas en efectivo para ``asegurar'' el capital.
\end{itemize}

\subsection{La Regla del Juego: La Función de Recompensa}

Si las acciones son los movimientos, la recompensa ($R$) es el sistema de ``premios y castigos''. No queremos que nuestra IA sea simplemente una apostadora; queremos que sea una inversora inteligente y equilibrada. Por eso, su sistema de evaluación tiene tres componentes:

\begin{equation}
    R_{t+1} = I_{t+1}r_{t+1} - c|I_{t+1} - I_t| - \lambda \sigma_t
\end{equation}

Esta fórmula es la ``calculadora de sabiduría'' de la IA y funciona así:
\begin{enumerate}
    \item \textbf{El Beneficio Directo ($I_{t+1}r_{t+1}$):} Es la ``zanahoria''. Si la IA decide estar comprada ($I=1$) y el precio sube, recibe un premio gigante. Si el precio baja, recibe un castigo.
    \item \textbf{El Impuesto al Movimiento ($-c|\dots|$):} Son las \textbf{comisiones}\footnote{Costo porcentual por cada operación ejecutada en el exchange; previene el exceso de operaciones innecesarias.} Cada vez que la IA compra o vende, Binance se queda con una pequeña parte. Esto castiga a la IA si intenta operar demasiado seguido sin una buena razón.
    \item \textbf{La Multa por el Riesgo ($-\lambda \sigma_t$):} Es el freno de seguridad. Si el mercado está muy inestable (volatilidad alta), la IA recibe una pequeña penalización. Esto le enseña a preferir la calma y a no arriesgarse innecesariamente en medio de tormentas financieras.
\end{enumerate}

\section{El Aprendizaje: ¿Cómo se vuelve inteligente nuestra IA?}

Una vez que tenemos el mundo definido (el mercado), el cuerpo (la IA) y el sistema de puntaje (la recompensa), falta la pieza más importante: el mecanismo de aprendizaje. En este proyecto, el cerebro de nuestra IA utiliza un algoritmo llamado \textbf{Q-Learning}\footnote{Algoritmo que aprende el valor de una acción en un estado dado mediante la experiencia directa.}.

\subsection{La ``Tabla de Sabiduría'' (Q-Table)}

Imagína que la IA lleva consigo un cuaderno gigante donde anota todo lo que aprende. Técnicamente, este cuaderno es una matriz matemática llamada \textbf{Q-Table}\footnote{Estructura de datos donde el agente asocia cada estado con el éxito esperado de cada acción posible.}. 

\begin{itemize}
    \item \textbf{Las Filas:} Representan todas las situaciones posibles (Estados). Por ejemplo: ``Precio bajando + Mucho miedo + Tengo dólares''.
    \item \textbf{Las Columnas:} Representan las decisiones posibles (Acciones): Comprar, Vender o Esperar.
    \item \textbf{Las Celdas (Valores Q):} En cada cruce de fila y columna, hay un número que representa qué tan buena idea es tomar esa acción en ese momento.
\end{itemize}

Al principio, la tabla está llena de ceros; la IA es un ``bebé'' que no sabe nada. Con cada hora que pasa en la simulación, la IA prueba una acción, recibe una recompensa y actualiza el número en su cuaderno mediante la \textbf{Ecuación de Bellman}\footnote{Base matemática que permite calcular el valor actual de una acción considerando premios futuros esperados.}:

\begin{equation}
    Q_{new}(s, a) = Q(s, a) + \alpha \underbrace{[R + \gamma \max Q(s', a') - Q(s, a)]}_{\text{Error de Aprendizaje}}
\end{equation}

Esta fórmula es el motor de inteligencia de nuestro agente y se traduce así:
\begin{itemize}
    \item \textbf{$Q(s, a)$:} Es lo que la IA sabía hasta ayer sobre esa situación.
    \item \textbf{$\alpha$ (Tasa de Aprendizaje}\footnote{Ritmo al que el agente incorpora nueva información, ponderando la experiencia reciente.}\textbf{):} Es el nivel de ``atención''. Si es alto, la IA olvida rápido lo viejo para creerle a lo nuevo. Si es bajo, aprende más despacio pero de forma más sólida.
    \item \textbf{$R$:} Es el premio o castigo inmediato que acaba de recibir.
    \item \textbf{$\gamma \max Q(s', a')$ (Visión de Futuro):} Es la ganancia máxima que la IA espera obtener en la siguiente hora. Esto es lo que permite que sea estratégica y no solo impulsiva.
    \item \textbf{El Corchete (Error):} Representa la ``sorpresa''. Es la diferencia entre lo que la IA esperaba ganar y lo que realmente pasó. Si la sorpresa es positiva, el valor Q sube; si es negativa, baja.
\end{itemize}

\subsection{El Dilema del Explorador: $\epsilon$-greedy}

Para que el aprendizaje sea efectivo, la IA se enfrenta a un dilema constante que nosotros también vivimos: ¿Hago lo que ya sé que funciona o pruebo algo nuevo por si acaso es mejor? Esto se conoce como el equilibrio entre Explotación y Exploración.

Utilizamos una estrategia llamada \textbf{Epsilon-greedy}\footnote{Estrategia que equilibra la explotación de lo aprendido con la exploración aleatoria de nuevas acciones.} ($\epsilon$-greedy), que funciona como un interruptor de curiosidad:

\begin{enumerate}
    \item \textbf{Explotación ($1-\epsilon$):} La mayor parte del tiempo (por ejemplo, el 90\% de las veces), la IA es ``obediente''. Mira su tabla y elige la acción que tiene el puntaje más alto. Aquí está usando su sabiduría acumulada para ganar dinero.
    \item \textbf{Exploración ($\epsilon$):} El resto del tiempo (el 10\% restante), la IA se vuelve ``exploradora''. Ignora su tabla y toma una decisión totalmente al azar. 
\end{enumerate}

\textbf{¿Para qué sirve el azar?} Sin este 10\% de curiosidad, si la IA encontrara una estrategia que gana un poquito de dinero, se quedaría haciendo eso para siempre y nunca descubriría una estrategia mucho mejor que requiere arriesgarse un poco más. El azar es lo que permite que nuestra IA pase de ser un trader mediocre a un inversor brillante.

\subsection{¿Por qué usar la Ecuación de Bellman?}
Muchos se preguntan por qué no evaluar las acciones solo por su resultado inmediato. La respuesta es simple: el trading es un juego de largo plazo. La \textbf{Ecuación de Bellman} es vital porque permite que el agente sea \textbf{estratégico}. Al incluir el valor futuro ($\gamma \max Q$), obligamos a la IA a entender que a veces es mejor perder un poco hoy para posicionarse y ganar mucho mañana. Es la diferencia entre un jugador impulsivo y un gran maestro de ajedrez.

\subsection{El porqué de Q-Learning Tabular (vs. Deep Learning)}
En la actualidad, muchos asocian la Inteligencia Artificial exclusivamente con Redes Neuronales Profundas (Deep Learning). Sin embargo, para este proyecto hemos elegido \textbf{Tabular Q-Learning} por tres razones estratégicas:
\begin{enumerate}
    \item \textbf{Transparencia e Interpretabilidad (Cero Caja Negra):} En finanzas, es crucial saber \textit{por qué} la IA tomó una decisión. Con una tabla Q, podemos abrir el ``cuaderno de sabiduría'' del agente y ver el valor numérico exacto asignado a cada decisión. En el Deep Learning, esto suele ser un misterio de ``caja negra''.
    \item \textbf{Eficiencia de Datos en Series de Tiempo:} El mercado cripto es extremadamente ruidoso. El Q-Learning tabular, apoyado en la discretización de estados, ignora las variaciones insignificantes y se enfoca en patrones lógicos de forma mucho más rápida y precisa que un modelo de aprendizaje profundo.
    \item \textbf{Control y Garantía de Convergencia:} El Q-Learning tabular garantiza matemáticamente que, bajo las condiciones adecuadas de exploración, el agente encontrará la política óptima. Esto nos da un respaldo teórico sólido que es mucho más errático de asegurar con redes neuronales en entornos tan volátiles.
\end{enumerate}

\section{El Factor Humano: Teoría de las Finanzas Conductuales}

¿Por qué incluir un índice de ``emociones'' en un proyecto de matemáticas y programación? La respuesta está en una rama moderna de la economía llamada \textbf{Finanzas Conductuales}.

\subsection{La Racionalidad Limitada}

Tradicionalmente, se creía que los inversores eran máquinas perfectas que siempre tomaban la decisión más lógica. Sin embargo, la realidad nos dice lo contrario: los humanos sufrimos de \textbf{sesgos}\footnote{Patrones de pensamiento irracional que afectan el juicio y la toma de decisiones.}. En el mercado cripto, esto se traduce en dos grandes fuerzas:
\begin{itemize}
    \item \textbf{FOMO (\textit{Fear Of Missing Out}):} El miedo a quedarse fuera de una subida, que empuja a la gente a comprar cuando los precios ya están muy caros (codicia).
    \item \textbf{Venta por Pánico:} El miedo que obliga a vender cuando los precios bajan, justo cuando el mercado está por recuperarse.
\end{itemize}

\subsection{El Índice como Puente}

El Crypto Fear \& Greed actúa como nuestro sensor de estos comportamientos irracionales. Al darle esta información a nuestra IA, no solo le estamos dando números, le estamos dando el contexto psicológico del mercado. Esto le permite aprender, por ejemplo, que si hay ``Miedo Extremo'', es muy probable que los precios estén artificialmente bajos y sea una oportunidad de oro.

\section{Criterios de Éxito: ¿Cómo evaluaremos a nuestro Agente?}

En el mundo real, no basta con decir que la IA ``ganó dinero''. Un mono lanzando dardos podría ganar dinero si el mercado sube por suerte. Para saber si nuestra IA es realmente inteligente, utilizaremos tres métricas de oro:

\subsection{Rentabilidad vs. Riesgo (Ratio de Sharpe)}

El \textbf{Ratio de Sharpe}\footnote{Métrica que compara la rentabilidad obtenida frente al riesgo asumido para verificar la eficiencia del sistema.} es la métrica más importante en las finanzas. Nos dice cuánto beneficio extra genera la IA por cada unidad de ``susto'' o riesgo que nos hace pasar. 
\begin{itemize}
    \item Un Sharpe alto significa que la IA es eficiente y gana dinero de forma tranquila y constante.
    \item Un Sharpe bajo significa que, aunque gane dinero, lo hace de forma muy arriesgada e inestable.
\end{itemize}

\subsection{El Máximo Dolor: Drawdown}

El \textbf{Máximo Drawdown}\footnote{Caída acumulada más profunda desde un punto máximo de la inversión; mide el riesgo de pérdida de capital.} mide la caída más grande que ha sufrido el capital desde su punto más alto. Es la medida del ``peor momento'': si la IA lleva mi cuenta de \$100 a \$150 y luego cae a \$80, el drawdown fue enorme. Nuestro objetivo es que la IA aprenda a proteger el capital para que estos baches sean lo más pequeños posible.

\subsection{Entrenamiento vs. Prueba (Generalización)}

Para asegurar que la IA no haya ``memorizado'' el pasado (algo que llamamos \textbf{overfitting}\footnote{Error de aprendizaje donde el modelo memoriza datos específicos del pasado en lugar de entender principios generales.}), dividiremos nuestros datos. La IA aprenderá en un periodo y luego la pondremos a prueba en un periodo que nunca ha visto. Si sobrevive y gana en el ``mundo desconocido'', sabremos que hemos creado un sistema robusto y no solo una calculadora de historia.
\end{document}
